% Procedencia: integral 2249, §25.10--13 y cap. 89; propietarios 02_pantalla_accion,
% 37_cuerdas_branas_teoria_m, terna_espectral_teoria_m, capitulo_25_pantallas_teoria_m.
\section{Acción sobre trayectorias y realización dinámica undecadimensional}
\label{iv:accion-realizacion}

La extensión de la estructura a una dinámica dimensional requiere
transportar conjuntamente el estado, los canales y los datos de
frontera. Esta sección comienza por un funcional definido sobre
las trayectorias del TPK y por su ley exacta de composición.
La superficie de mundo y el codominio undecadimensional se
introducen después, con los mapas que deben conservar esa
composición y los operadores ya construidos.

\subsection{Descenso de una evolución y suficiencia del dato observado}

Sea \(q:\mathcal X_{\mathrm{HMT}}\to Y\) una lectura
sobreyectiva de los estados con memoria, y sea
\(U:\mathcal X_{\mathrm{HMT}}\to\mathcal X_{\mathrm{HMT}}\)
una evolución completa. Una evolución sobre los valores
observados existe precisamente cuando el dato conservado por
\(q\) basta para determinar el siguiente valor observado.

\begin{theorem}[Criterio exacto de descenso por fibras]
\label{iv:thm:descenso-fibras}
Existe una única aplicación \(\overline U:Y\to Y\)
que satisfaga \(qU=\overline Uq\) si y sólo si
\begin{equation}
 q(x)=q(y)\quad\Longrightarrow\quad q(Ux)=q(Uy).
 \label{iv:eq:descenso-fibras}
\end{equation}
\end{theorem}
\begin{proof}
Si existe \(\overline U\), aplicar la misma función a
\(q(x)=q(y)\) prueba la condición. Recíprocamente, se
define \(\overline U(q(x))=q(Ux)\). La condición
\eqref{iv:eq:descenso-fibras} garantiza independencia del
representante. La sobreyectividad de \(q\) define la
aplicación en todo \(Y\) y prueba su unicidad.
\end{proof}

El teorema permite decidir qué memoria debe conservar una
lectura. Cuando dos representantes tienen imágenes observadas
diferentes, se puede ampliar el dato retenido, restringir
el dominio o utilizar una correspondencia multivaluada.
Son elecciones matemáticas distintas. En particular, el
retorno de una coordenada de fase no autoriza a identificar
las evoluciones de dos historias que conserven memorias
diferentes.

\subsection{Cociclo de cambios direccionales y aritméticos}

Una trayectoria admisible del TPK lleva, en cada arista
\(e_j\), una dirección \(d_j\) y un tipo aritmético
\(t_j\). Para contabilizar los cambios en una concatenación,
se conserva también el dato precedente. La arista aumentada es
\(\widetilde e_j=(e_j,d_{j-1},t_{j-1})\); la pareja
\((d_0,t_0)\) forma parte de la frontera inicial. Definimos
\[
 w_{\mathrm{giro}}(\widetilde e_j)
       =\mathbf1_{\{d_j\ne d_{j-1}\}},\qquad
 w_{\mathrm{tipo}}(\widetilde e_j)
       =\mathbf1_{\{t_j\ne t_{j-1}\}}.
\]
Para \(\widetilde\gamma=(\widetilde e_1,\ldots,
\widetilde e_R)\), sean
\begin{equation}
 \mathbf I(\widetilde\gamma)=
 \left(\sum_{j=1}^R w_{\mathrm{giro}}(\widetilde e_j),
       \sum_{j=1}^R w_{\mathrm{tipo}}(\widetilde e_j)\right)
 =\bigl(N_{\mathrm{giro}},N_{\mathrm{tipo}}\bigr).
 \label{iv:eq:cociclo-accion}
\end{equation}
Dos trayectorias se componen cuando coinciden sus extremos
y el dato precedente de la segunda coincide con el dato
terminal de la primera.

\begin{proposition}[Aditividad con frontera conservada]
\label{iv:prop:accion-aditiva}
Para toda pareja composable,
\begin{equation}
 \mathbf I(\widetilde\gamma\widetilde\eta)
 =\mathbf I(\widetilde\gamma)+\mathbf I(\widetilde\eta).
 \label{iv:eq:accion-aditiva}
\end{equation}
Por tanto, \(\mathbf I\) es un cociclo aditivo con valores
en \(\mathbb N^2\) sobre la categoría de trayectorias
aumentadas. Para cada \(\lambda>0\), el funcional
\[
 I_\lambda(\widetilde\gamma)
 =N_{\mathrm{giro}}(\widetilde\gamma)
      +\lambda N_{\mathrm{tipo}}(\widetilde\gamma)
\]
es también aditivo.
\end{proposition}
\begin{proof}
La concatenación no cambia el dato precedente de ninguna
arista: el de la primera arista de \(\widetilde\eta\)
coincide, por la condición de composición, con el terminal
de \(\widetilde\gamma\). La suma sobre todas las aristas
se divide exactamente en las dos sumas de
\eqref{iv:eq:accion-aditiva}. Aplicar el funcional lineal
\((a,b)\mapsto a+\lambda b\) prueba la última afirmación.
\end{proof}

La información de frontera es parte efectiva de la prueba.
Si dos segmentos de una arista se evalúan por separado
ignorando el enlace entre ellos, ambos pueden tener cuenta
de giro nula aunque sus direcciones sean diferentes. El
camino concatenado contiene entonces un giro. Conservar
\((d_{j-1},t_{j-1})\) registra ese término en la primera
arista del segundo segmento y restituye la aditividad.

Para extremos y datos de frontera fijados, el principio de
interacción mínima del corpus selecciona minimizadores de
\(I_\lambda\) entre las rutas admisibles. El principio
es una regla de selección, mientras que el siguiente
resultado establece una propiedad de existencia finita.

\begin{proposition}[Existencia finita de minimizadores]
Si el conjunto \(\Gamma(a,b)\) de trayectorias admisibles
con la frontera fijada es finito y no vacío, existe
\(\gamma_*\in\Gamma(a,b)\) tal que
\(I_\lambda(\gamma_*)\leq I_\lambda(\gamma)\)
para todo \(\gamma\in\Gamma(a,b)\).
\end{proposition}
\begin{proof}
La imagen de \(\Gamma(a,b)\) por \(I_\lambda\) es
un conjunto finito no vacío de números reales, que posee
un elemento mínimo y al menos una preimagen.
\end{proof}

La proposición no selecciona el valor de \(\lambda\),
ni confunde existencia con unicidad. El valor y la regla
de selección utilizados en una realización determinada
deben formar parte de sus datos declarados. El funcional
aditivo recuperado aquí cuantifica cambios del propio
camino; su transporte a una acción dimensional constituye
una operación ulterior.

\subsection{Persistencia extendida y superficie de mundo}

La realización de una ruta como objeto extendido distribuye
su genealogía sobre una coordenada interna. Para una cuerda,
la historia tiene dos coordenadas locales \(\sigma^a\):
una espacial interna y otra temporal. Una realización se
describe por un mapa \(X:\Sigma\to\mathcal M\), una
métrica \(h_{ab}\) en la superficie de mundo y campos
de fondo en el espacio objetivo. Las condiciones sobre
\(\partial\Sigma\) forman parte del dominio variacional.
Para una brana de dimensión espacial \(p\), la historia
es un volumen de mundo de dimensión \(p+1\).

El codominio variacional de cuerda que utiliza el corpus
es la acción de Polyakov
\begin{equation}
 S_{\mathrm P}[X,h]
 =-\frac{T_{\mathrm s}}2\int_\Sigma d^2\sigma\,
 \sqrt{-h}\,h^{ab}g_{MN}(X)
             \partial_aX^M\partial_bX^N,
 \label{iv:eq:polyakov}
\end{equation}
donde \(T_{\mathrm s}\) es la tensión. La notación
\(T_{\mathrm s}\) la distingue de la dualidad
\(\mathsf T_0\), del transporte cardinal \(T_d\)
y de la evolución completa \(U\).

La composición de trayectorias aumentadas tiene ya la
propiedad demostrada en \ref{iv:prop:accion-aditiva}.
Una realización continua de esa composición debe
especificar el mapa de caminos a configuraciones,
su transporte de las condiciones de frontera y la
normalización que compara \(I_\lambda\) con la acción
continua. La integral \eqref{iv:eq:polyakov} precisa
el codominio y sus campos; sus parámetros no eligen
retroactivamente las emisiones ni los registros del TPK.

\subsection{Realización undecadimensional y operadores graduados}

La reducción dirigida por \(K\) proporciona
\(V_{11}=\ell_K\oplus V_{10}\), y el teorema
\ref{iv:thm:pantallas-isomorfas} coordina esa reducción
con la pantalla reticular y la dualidad. La realización
física se tipa mediante
\[
 \Xi_M:\mathscr S_M^{\mathrm{alg}}
                  \longrightarrow\mathcal X_{11}^{\mathrm{phys}}.
\]
El codominio lleva métrica \(g_{MN}\), tres-forma
\(C_{MNP}\), gravitino \(\psi_M\), acción,
restricciones de calibre y transformaciones de
supersimetría. Una realización operatoria incluye
\[
 \mathfrak S_M=(\mathcal H_{\mathrm{phys}},H,Q),\qquad
 \mathcal H_{\mathrm{phys}}=\mathcal H_+\oplus\mathcal H_-.
\]
La graduación se expresa mediante una involución
\(\Gamma\), igual a \(+I\) en \(\mathcal H_+\)
y a \(-I\) en \(\mathcal H_-\). La supercarga
es impar, \(\Gamma Q=-Q\Gamma\), y el
Hamiltoniano es no negativo. La relación
\(Q^2=H\) se entiende como igualdad de operadores,
incluyendo sus dominios.

Sea \(\mathcal H_{\mathrm{HMT}}\) una realización
hilbertiana de los estados genealógicos pertinentes,
con operadores \(\widehat H,\widehat Q\). Una
isometría \(W:\mathcal H_{\mathrm{HMT}}\to
\mathcal H_{\mathrm{phys}}\) debe transportar
dominios, graduación y observables. Las condiciones
operatorias se escriben
\begin{equation}
 W\widehat H=HW,\qquad
 W\widehat Q=QW,\qquad
 WU_T^{\mathrm{HMT}}=U_T^{\mathrm{phys}}W.
 \label{iv:eq:entrelazamiento-fisico}
\end{equation}

\begin{theorem}[Transporte compatible de la supercarga]
\label{iv:thm:supercarga-transporte}
Supóngase \(\widehat Q^{\,2}=\widehat H\), que
\(W D(\widehat Q)\subset D(Q)\), y que los
entrelazamientos de \eqref{iv:eq:entrelazamiento-fisico}
se satisfacen en sus dominios. Entonces, para
\(\psi\in D(\widehat H)=D(\widehat Q^{\,2})\),
\[
 W\psi\in D(Q^2),\qquad Q^2W\psi=HW\psi.
\]
Si además el transporte físico de los radios respeta
\eqref{iv:eq:conjugacion-h}, el sector compacto
transportado conserva su espectro bajo radios recíprocos.
\end{theorem}
\begin{proof}
Como \(\psi\in D(\widehat Q^2)\), tanto
\(\psi\) como \(\widehat Q\psi\) pertenecen
a \(D(\widehat Q)\). Por tanto, \(W\psi\in D(Q)\)
y \(QW\psi=W\widehat Q\psi\in D(Q)\). Se puede
aplicar \(Q\) de nuevo, y
\[
 Q^2W\psi=W\widehat Q^{\,2}\psi
          =W\widehat H\psi=HW\psi.
\]
La igualdad espectral del sector compacto se transporta
por la isometría y el entrelazamiento unitario de T;
sobre la imagen cerrada de \(W\), éste es una
equivalencia unitaria de las realizaciones correspondientes.
\end{proof}

El teorema conserva explícitamente los operadores y el
dominio sobre el que se prueba la identidad. La
graduación organiza dos sectores; la construcción de
campos fermiónicos y sus relaciones de ocupación posee
su propia estructura. Las distribuciones térmicas y el
catálogo de especies no se utilizan en la demostración
de las pantallas ni en la dualidad compacta de este
artículo.

\subsection{Compactificación compatible y conclusión del desarrollo}

Una compactificación del estado se describe mediante
una sección prolongable \(x_\infty\), un mapa
\(\kappa_{\mathrm{comp}}\) hacia el espacio de módulos
elegido y un transporte \(\mathcal U\) de los datos
de métrica, orientación, frontera y memoria. El
subíndice distingue este mapa de la lectura racional
\(\kappa\) del registro dodecafásico. Su compatibilidad
exige cuatro propiedades concretas:
\begin{enumerate}
\item las restricciones finitas de \(x_\infty\)
      pertenecen a las prolongaciones admisibles del TPK;
\item el mapa de compactificación entrelaza las dualidades,
      \(\kappa_{\mathrm{comp}}\mathsf T
        =\mathsf T_M\kappa_{\mathrm{comp}}\);
\item los observables y el espectro descienden a los
      cocientes elegidos conforme al criterio
      \ref{iv:thm:descenso-fibras};
\item la realización transporta la dirección
      \(\ell_K\), las condiciones de frontera y los
      dominios operatorios, y conserva las ecuaciones
      de acción y supersimetría declaradas.
\end{enumerate}

Las pruebas precedentes proporcionan la realización
hilbertiana de fase, el cierre tracial, las pantallas
con sus mapas de cociente, la involución de radio,
su equivalencia unitaria y el cociclo aditivo de
trayectorias. La realización undecadimensional
especifica los campos y las compatibilidades que
utiliza al recibir esa estructura. Su identificación
física completa conserva, además, las condiciones
de tensión, osciladores, amplitudes, anomalías y
límite de baja energía que enumera el corpus.
Estas condiciones no alteran las identidades
algebraicas demostradas: delimitan la afirmación
posterior que se obtiene al transportarlas.

La prolongación excepcional y la dimensional
permanecen vinculadas por su origen común. El
Monstruo actúa en la construcción graduada
\(V^\natural\); el Hamiltoniano compacto y las
pantallas reciben los datos de otros mapas del
mismo estado. Esta formulación conserva el
contenido positivo de la doble realización,
con sus operaciones, sus pruebas y la información
transmitida en cada etapa (teorema~\ref{iv:thm:supercarga-transporte}).
